\section{Worked Example: TLA+ to Code --- Finding and Fixing a Validation Bug}
\index{TLA+}\index{TLC}\index{counterexample}\index{validation bug}\index{NOrec}\index{model checking!trace to code}
\label{sec:tla-to-code}

The earlier worked examples showed a broken PlusCal model and the TLC
counterexample it produced.  This section completes the loop: the same
stale-snapshot bug is traced from the model to the actual C++ backend, the
corresponding code is identified, and the fix is verified by re-running TLC.

\subsection{The Broken Model}

Consider a simplified NOrec commit action that omits the value recheck:

\begin{lstlisting}[caption={A deliberately broken commit action: write-back without re-checking read-set values.}]
BrokenCommit ==
  /\ ok
  /\ ~commitLock
  /\ commitLock' = TRUE
  /\ IF r1 >= amount THEN
        /\ bal' = [bal EXCEPT ![1] = @ - amount,
                              ![2] = @ + amount]
        /\ clock' = clock + 1
     ELSE
        /\ bal' = bal
        /\ clock' = clock
\end{lstlisting}

The guard \texttt{ok} was set earlier by the optimistic validation:
if the clock had not changed since the snapshot, \texttt{ok} became
\texttt{TRUE}.  But the commit action above does not re-read
\texttt{bal[1]} and \texttt{bal[2]} under the lock before writing back.
It trusts the cached snapshot values \texttt{r1} and \texttt{r2}
unconditionally.

\subsection{The TLC Counterexample}

TLC finds a short two-transaction trace.  The schedule is:

\begin{enumerate}
  \item Transaction~1 records \texttt{snap = 0}, reads
        \texttt{r1 = bal[1] = 10} and \texttt{r2 = bal[2] = 10}.
  \item Transaction~2 transfers 5 from account~1 to account~2.
        \texttt{bal[1] = 5}, \texttt{bal[2] = 15}, \texttt{clock = 1}.
  \item Transaction~1's optimistic validation now fires: the clock has
        changed (\texttt{clock = 1}, \texttt{snap = 0}), so it
        re-reads both accounts, finds them equal to its cached
        \texttt{r1} and \texttt{r2}, and sets \texttt{ok = TRUE}.
  \item Transaction~1 enters the broken commit, which acquires the lock
        and writes back \texttt{bal[1] = 10 - 1 = 9},
        \texttt{bal[2] = 10 + 1 = 11}.
\end{enumerate}

The final state is \texttt{bal[1] = 9}, \texttt{bal[2] = 11} ---
total money is 20, but account~1 was debited using a stale snapshot
from \emph{before} Transaction~2's transfer.  The correct outcome
should be \texttt{bal[1] = 4}, \texttt{bal[2] = 16}.  TLC reports:

\begin{lstlisting}[caption={TLC counterexample output (simplified).}]
Invariant MoneyConserved violated.
State 4: bal = <<9, 11>>, clock = 1, committed = <<TRUE, TRUE>>
\end{lstlisting}

The invariant that fails is \texttt{MoneyConserved == Total(bal) = 20}:
the total is still 20, but the money distribution is wrong.  A sharper
invariant---one that tracks whether each individual transfer's debit
and credit are both visible---would also fail.

\subsection{Tracing to C++}

The corresponding C++ implementation is
\rep{backends/tm\_impl/norec/NOrec.hpp}.
The \texttt{commit()} function (line~251) implements the protocol that
the TLA+ model checks.  The critical validation logic lives inside the
CAS retry loop:

\begin{lstlisting}[caption={Simplified NOrec commit path (\texttt{NOrec.hpp:251}).}]
inline void commit() {
  auto *tx = current_tx;
  if (tx->read_only) { tx->reset(); return; }

  word_t expect = tx->snapshot;
  while (!global_lock.compare_exchange_strong(
             expect, expect + 1,
             std::memory_order_acquire,
             std::memory_order_relaxed)) {
    // CAS failed: clock changed.  Re-check values.
    tx->snapshot = validate();
    expect = tx->snapshot;
  }
  // Lock acquired.  Write back.
  for (auto &w : tx->write_set)
    write_value_to_addr(w.addr, w.new_val, w.type);
  set_clock(expect + 2);       // release lock, advance clock
  tx->reset();
}
\end{lstlisting}

The \texttt{validate()} function (line~227) is where the value recheck
happens:

\begin{lstlisting}[caption={The value validation that the broken model omitted.}]
inline word_t validate() {
  auto *tx = current_tx;
  volatile word_t time = 0;
  while (true) {
    time = get_clock();
    if ((time & 1) != 0) continue;     // lock held, spin
    for (auto &r : tx->read_set) {
      any_type_t check =
          read_value_from_addr(r.addr, r.type);
      if (!iseq(check, r.observed_val))  // THE CRITICAL CHECK
        abort_tx();                      // stale read -> abort
    }
    if (time == get_clock()) return time;  // clock stable
  }
}
\end{lstlisting}

If the line \texttt{if (!iseq(check, r.observed\_val)) abort\_tx();}
were removed, the function would always succeed: the read-set entries
would be walked but never actually checked.  A transaction that read
stale values would proceed to write-back and publish incorrect results
--- the same bug the TLA+ model exposed.

The connection between model and code is direct:

\begin{itemize}
  \item The PlusCal commit action's missing recheck corresponds to
        removing the \texttt{iseq()} comparison in
        \texttt{validate()}.
  \item The CAS retry in the C++ code corresponds to the
        \texttt{elsif clk = snapshot[self]} branch in the PlusCal
        model: if the clock changed, the transaction must re-validate
        before attempting the CAS again.
  \item The value-based comparison (\texttt{iseq}) is the same check
        that the correct PlusCal model encodes as
        \texttt{r1 = bal[1] /\ r2 = bal[2]}.
\end{itemize}

\subsection{The Fix}

The fix in the model is to restore the value recheck:

\begin{lstlisting}[caption={The correct commit guard.}]
CommitRecheck ==
  /\ clock = snap
  /\ r1 = bal[1]
  /\ r2 = bal[2]
\end{lstlisting}

In the C++ code, the recheck already exists as the \texttt{iseq()} test
inside \texttt{validate()}.  The lesson is that the model and the code
encode the same safety argument in different notation.  When someone
removes the recheck from the model, TLC catches the bug.  If someone
removes the corresponding line from \texttt{NOrec.hpp}, the bank
benchmark loses money --- the same invariant that TLC checked, violated
in the real implementation.

\subsection{Re-running TLC}

After restoring the recheck, TLC re-runs and confirms the fix:

\begin{lstlisting}[caption={TLC output after the fix is restored.}]
Model checking completed.  No errors found.
149430 states generated, 64891 distinct states.
\end{lstlisting}

The model and the code agree: the value recheck is the single line
that separates a correct NOrec from a stale-snapshot mutation.  This is
the end-to-end workflow of \S\ref{sec:verify-workflow}: specify, check,
inject a bug, confirm the checker finds it, trace the result to code,
fix, and re-check.
